% Đề cương chi tiết — bắt buộc theo checklist nộp cuốn báo cáo của Khoa
% (đặt giữa Lời cảm ơn và Mục lục). Bản nộp TRƯỚC bảo vệ không cần chữ ký;
% bản nộp SAU bảo vệ phải có chữ ký GVHD — khi đó thay toàn bộ nội dung
% file này bằng \includepdf bản scan đã ký.

\chapter*{Đề cương chi tiết}
\markboth{Đề cương chi tiết}{Đề cương chi tiết}

\section*{1.\quad THÔNG TIN CHUNG}

\begin{itemize}
  \item \textbf{Tên đề tài:} Xây dựng dữ liệu và tinh chỉnh mô hình ngôn ngữ
  nhỏ cho tác nhân trí tuệ nhân tạo trên web \\
  \textit{(Data Construction and Fine-tuning of Small Language Models for Web
  AI Agents)}
  \item \textbf{Người hướng dẫn:} PGS.TS. Nguyễn Văn Vũ (Khoa Công nghệ
  Thông tin, Trường Đại học Khoa học Tự nhiên, ĐHQG-HCM)
  \item \textbf{Sinh viên thực hiện:}
  \begin{enumerate}
    \item Trương Bảo Kiệt - MSSV 22127223
    \item Hồ Tiến Phát - MSSV 22127321
  \end{enumerate}
  \item \textbf{Loại đề tài:} Nghiên cứu
  \item \textbf{Thời gian thực hiện:} 01/2026 - 08/2026
\end{itemize}

\section*{2.\quad NỘI DUNG THỰC HIỆN}

\subsection*{2.1\quad Giới thiệu về đề tài}

\subsubsection*{2.1.1\quad Bối cảnh và vấn đề}

Trình duyệt là giao diện chung của hầu hết các dịch vụ trực tuyến, nên một
agent phần mềm thao tác được website \emph{qua chính giao diện người dùng} là
một trong những dạng tự động hóa có phạm vi ứng dụng rộng nhất. Các mô hình
ngôn ngữ lớn (LLM) gần đây đã khiến những web agent như vậy trở nên khả thi:
nhận một bản mô tả dạng văn bản của trang cùng một mục tiêu, mô hình quyết
định được cần click hay gõ gì tiếp theo.

Tuy nhiên, năng lực đó hiện gần như chỉ tồn tại ở các mô hình độc quyền rất
lớn truy cập qua API, và sự phụ thuộc này gây tốn kém theo nhiều cách cộng
dồn. Mỗi bước của mỗi task là một lời gọi API trả phí trong khi task của agent
cần rất nhiều bước; độ trễ bị giới hạn bởi dịch vụ từ xa; toàn bộ nội dung mọi
trang được duyệt (có thể chứa dữ liệu riêng tư hoặc nội bộ) bị gửi cho
bên thứ ba; và bản thân mô hình là một hàng hóa cố định mà người vận hành
không kiểm tra hay cải thiện được.

Mô hình ngôn ngữ nhỏ (SLM, cỡ mười tỷ tham số trở xuống, tự triển khai
được trên một GPU) đảo ngược mọi ràng buộc trên. Thứ chúng thiếu là năng
lực. Ngay cả các mô hình lớn cũng thấy task web thực tế là khó: trên benchmark
WebArena gồm các website đầy đủ chức năng tự triển khai, agent dựa trên GPT-4
chỉ hoàn thành khoảng 14\% số task so với khoảng 78\% của con
người~\cite{zhou2024webarena}. Phương thuốc tiêu chuẩn, tinh chỉnh có giám
sát (SFT) trên các bản mẫu hành vi, lại vấp phải bài toán dữ liệu:
trajectory tương tác có gán mục tiêu trên website đích hầu như không tồn tại,
còn thuê người gán nhãn thì đắt, chậm và phải làm lại cho từng website.

Câu hỏi đề tài đặt ra là: liệu có thể xây dựng một pipeline hoàn toàn tự động
\emph{khám phá} một website không nhãn, \emph{tự nghĩ ra} các task có ý nghĩa
dựa trên những gì nó tìm thấy, \emph{kiểm chứng} nỗ lực nào thực sự thành
công, rồi biên dịch kết quả thành dữ liệu huấn luyện cho một SLM? Nếu được,
chi phí chuyên hóa một SLM cho một website giảm từ một chiến dịch gán nhãn thủ
công xuống thành chi phí tính toán.

\subsubsection*{2.1.2\quad Các nghiên cứu liên quan}

\textbf{WebArena}~\cite{zhou2024webarena} xây dựng môi trường web tự triển
khai với 812 task trên năm loại website, chấm điểm bằng bộ đánh giá tự động
dựa trên trạng thái cuối thay vì khớp hành động. Đề tài dùng WebArena làm cả
môi trường khám phá lẫn benchmark đánh giá, và dùng
BrowserGym~\cite{dechezelles2025browsergym} làm lớp giao tiếp trình duyệt.

\textbf{Go-Browse}~\cite{gandhi2026gobrowse} là công trình gần đề tài nhất:
khám phá website có cấu trúc bằng một graph các trang để sinh dữ liệu huấn
luyện, đạt 21,7\% trên WebArena với mô hình 7B. Tuy nhiên graph của Go-Browse
chỉ ghi \emph{trang}; nó không trừu tượng hóa các năng lực tái sử dụng, nên
task sinh ra vẫn bám phạm vi một trang. Đề tài phân biệt bằng graph
\emph{hai tầng} tách biệt "agent đã ở đâu" khỏi "agent làm được gì", cho phép
ghép chuỗi năng lực thành task dài hạn.

\textbf{SkillWeaver}~\cite{zheng2025skillweaver} và
\textbf{ASI}~\cite{wang2025asi} phát hiện và tích lũy skill tái sử dụng trong
quá trình agent hoạt động. Đề tài kế thừa ý tưởng skill nhưng nâng skill đã
kiểm chứng thành \emph{operator} có precondition/effect tường minh, để có thể
lập kế hoạch bằng tìm kiếm trên graph.

\textbf{NNetNav}~\cite{murty2024nnetnav}, \textbf{OS-Genesis}~\cite{sun2024osgenesis}
và \textbf{AgentTrek}~\cite{xu2024agenttrek} sinh dữ liệu huấn luyện agent tự
động bằng cách tương tác rồi gán nhãn ngược lại mục tiêu. Đề tài khác ở chỗ
mục tiêu không được suy ngược từ một chuỗi hành động bất kỳ mà được
\emph{ghép có chủ đích} từ các năng lực đã kiểm chứng.

\textbf{LoRA}~\cite{hu2022lora} là kỹ thuật tinh chỉnh hiệu quả tham số, chèn
các ma trận phân rã hạng thấp vào lớp Attention thay vì cập nhật toàn bộ trọng
số. Đề tài áp dụng LoRA để fine-tune mô hình nền
Qwen3.5-9B~\cite{yang2025qwen3}, huấn luyện qua
LLaMA-Factory~\cite{zheng2024llamafactory} và phục vụ suy luận bằng
vLLM~\cite{kwon2023vllm}.

Về mặt hình thức hóa, biểu diễn skill của đề tài dựa trên khung
\emph{options} trong học tăng cường~\cite{sutton1999options}, còn biểu diễn
operator dựa trên khung precondition/effect kiểu STRIPS~\cite{fikes1971strips}.
Việc chuẩn hóa mệnh đề dùng sentence embedding~\cite{reimers2019sbert}.

\subsubsection*{2.1.3\quad Hệ thống thực tế khảo sát}

Đề tài khảo sát và thực nghiệm trên họ website tự triển khai của benchmark
WebArena, các bản sao đầy đủ chức năng của phần mềm thật, có dữ liệu thật
và trạng thái thay đổi được:

\begin{itemize}
  \item \textbf{Reddit} (diễn đàn): forum, bài đăng, bình luận, bỏ phiếu,
  trang cá nhân.
  \item \textbf{GitLab} (kho mã nguồn): dự án, issue, merge request, thành
  viên, cài đặt.
  \item \textbf{Shopping} (thương mại điện tử): danh mục, chi tiết sản phẩm,
  giỏ hàng, thanh toán.
  \item \textbf{Shopping Admin} (trang quản trị): đơn hàng, tồn kho, khách
  hàng, báo cáo.
\end{itemize}

\subsection*{2.2\quad Mục tiêu đề tài}

\textbf{Vấn đề nghiên cứu.}

\begin{quote}
Cho (i) một website đang hoạt động, không có bất kỳ nhãn task nào, truy cập
qua trình duyệt tiêu chuẩn, và (ii) một mô hình ngôn ngữ nhỏ đã pretrain,
hãy tạo ra phiên bản fine-tune của mô hình đó có khả năng hoàn thành các
task phát biểu bằng ngôn ngữ tự nhiên trên website, mà không dùng bất kỳ
bản mẫu hay nhãn task nào do con người tạo.
\end{quote}

\textbf{Các mục tiêu cụ thể.} Đề tài theo đuổi năm mục tiêu, tương ứng một--một
với năm thách thức của bài toán:

\begin{itemize}
  \item[\textbf{O1.}] Xây dựng pipeline khám phá tự động thu thập trajectory
  tương tác từ một website không nhãn, bền vững trước cấu trúc trùng lặp và
  trang lỗi, an toàn trên trạng thái dùng chung.
  \item[\textbf{O2.}] Sinh các \emph{semantic task dài hạn đã kiểm chứng} bằng
  cách phát hiện skill tái sử dụng, nâng chúng thành operator lập kế hoạch, và
  ghép chuỗi các operator trên graph.
  \item[\textbf{O3.}] Tinh lọc trajectory thu được thành tập dữ liệu
  fine-tuning qua các cổng kiểm soát chất lượng tự động, dưới nguyên tắc
  \emph{prompt parity} nghiêm ngặt giữa huấn luyện và vận hành.
  \item[\textbf{O4.}] Fine-tune một mô hình ngôn ngữ nhỏ trên tập dữ liệu thu
  được và đánh giá trên bộ task chính thức của bốn website WebArena.
  \item[\textbf{O5.}] Đúc kết kinh nghiệm kỹ thuật khi vận hành LLM agent trên
  website thật thành các bài học và best practice tái sử dụng được.
\end{itemize}

Mã nguồn pipeline, tập dữ liệu thu thập và tài liệu tái lập thực nghiệm dự
kiến được công bố sau khi khóa luận hoàn thành.

\subsection*{2.3\quad Phạm vi của đề tài}

\subsubsection*{2.3.1\quad Các giải pháp được thử nghiệm và so sánh}

\begin{itemize}
  \item \textbf{Baseline:} mô hình nền Qwen3.5-9B-Instruct chưa fine-tune,
  đánh giá trên bộ task chính thức của WebArena; đây là ngưỡng tham chiếu
  cần vượt qua.
  \item \textbf{Giải pháp đề xuất:} chính mô hình đó sau khi fine-tune bằng
  LoRA trên tập dữ liệu do pipeline HSGE tự sinh.
  \item \textbf{So sánh nội bộ (ablation):} dữ liệu khám phá theo graph so với
  khám phá phẳng cùng ngân sách; có và không có cổng judge; có và không có
  bước viết lại suy nghĩ.
\end{itemize}

\subsubsection*{2.3.2\quad Mã nguồn và dữ liệu sẵn có}

\begin{itemize}
  \item \textbf{WebArena}\footnote{\url{https://github.com/web-arena-x/webarena}}:
  website, bộ task và bộ chấm điểm tự động; dùng nguyên trạng làm môi trường
  và benchmark.
  \item \textbf{BrowserGym}\footnote{\url{https://github.com/ServiceNow/BrowserGym}}:
  lớp giao tiếp trình duyệt (accessibility tree, action space, ảnh chụp màn
  hình) trên nền Playwright.
  \item \textbf{Qwen3.5-9B}\footnote{\url{https://huggingface.co/Qwen}}: mô
  hình nền được fine-tune.
  \item \textbf{LLaMA-Factory, vLLM, sentence-transformers}: lần lượt dùng cho
  huấn luyện LoRA, phục vụ suy luận, và mã hóa embedding phục vụ gom cụm skill
  và chuẩn hóa mệnh đề.
\end{itemize}

\subsubsection*{2.3.3\quad Kết quả thực nghiệm đã đạt được}

Tính đến thời điểm lập đề cương này, pipeline đã được cài đặt trọn vẹn và vận
hành trên website thật. Các con số dưới đây là hiện trạng tại thời điểm nộp đề
cương; kết quả cuối cùng của khóa luận được báo cáo ở
Chương~\ref{ch:experiments}:

\begin{itemize}
  \item \textbf{Khám phá:} 12 lượt chạy trên hai instance (Reddit và GitLab),
  thăm 301 node, ghi 1.174 episode (827 task đạt, 330 hỏng, 48 được agent
  tuyên bố bất khả thi).
  \item \textbf{Phát hiện skill:} 37 skill được kiểm chứng và thu nhận, nâng
  thành 35 operator lập kế hoạch có precondition/effect tường minh.
  \item \textbf{Tinh lọc dữ liệu:} tập dữ liệu SFT gồm 370 mẫu (352 huấn
  luyện / 18 kiểm định), dựng dưới ràng buộc prompt parity với agent lúc chạy
  thật.
  \item \textbf{Fine-tuning:} Qwen3.5-9B + LoRA huấn luyện thành công trên GPU
  A100-80GB thuê ngoài (loss huấn luyện 0,32; loss kiểm định 0,24).
  \item \textbf{Đánh giá sơ bộ:} trên bộ task chính thức của GitLab, mô hình
  nền đạt 10/180 và mô hình fine-tune đạt 12/180. Mỗi task mới chỉ được chạy
  \emph{một lần}, nên con số này là tín hiệu định hướng chứ chưa đủ cơ sở
  thống kê để kết luận.
\end{itemize}

Một hạn chế cần nêu rõ: \emph{chưa có semantic task ghép chuỗi nào được thực
thi} trong cả 12 lượt chạy. Thủ tục ghép chuỗi operator đã được cài đặt và có
kiểm thử tất định trên backend mô phỏng, nhưng trên site thật operator graph
gần như không có cạnh, do bước chuẩn hóa mệnh đề bằng độ tương tự embedding
không tách được synonym thật khỏi các cụm từ không liên quan, nên điều kiện
kích hoạt không xảy ra. Khắc phục hạn chế này thuộc phần công việc còn lại
(mục~2.6).

\subsubsection*{2.3.4\quad Những gì không thuộc phạm vi đề tài}

\begin{itemize}
  \item Phương pháp huấn luyện chỉ gồm tinh chỉnh có giám sát với LoRA; học
  tăng cường và tối ưu theo sở thích (DPO) nằm ngoài phạm vi.
  \item Chỉ thực nghiệm trên các instance WebArena tự triển khai; không tương
  tác với website công cộng ngoài benchmark.
  \item Không nghiên cứu tự động hóa cho giao diện native trên di động
  (Android/iOS).
  \item Không giải CAPTCHA hay né tránh cơ chế chống bot.
  \item Không xây dựng hệ thống RPA thương mại đầy đủ, cũng không triển khai
  mô hình như một dịch vụ API công khai; đây là công trình nghiên cứu.
  \item Toàn bộ task bằng tiếng Anh; đa ngữ nằm ngoài phạm vi.
\end{itemize}

\subsection*{2.4\quad Cách tiếp cận dự kiến}

\subsubsection*{2.4.1\quad Kiến trúc tổng thể}

Đề tài đề xuất \textbf{HSGE} (\textit{Hierarchical Skill-Graph Exploration})
một pipeline khám phá website dựa trên graph, vận hành như một vòng lặp
khép kín. Một \emph{explorer} đề xuất các task cụ thể bám trang trên mỗi trang
được thăm; một agent \emph{solver} thử thực hiện chúng trong trình duyệt thật,
sinh ra trajectory. Thành phần phát hiện cắt trajectory thành đoạn, gom cụm
các đoạn lặp lại, trừu tượng hóa mỗi cụm đủ chín thành một \emph{skill} tái sử
dụng, rồi nâng skill đã kiểm chứng thành \emph{operator} lập kế hoạch. Ghép
chuỗi operator trên graph sinh ra các task dài hạn, gọi là \emph{semantic task},
được đưa ngược lại chính solver đó, khép kín vòng lặp.

Ba nguyên tắc thiết kế xuyên suốt: (P1) \emph{kiểm chứng trước khi tái sử
dụng}: không gì được tin ngay lần đầu; (P2) \emph{graph thay vì phẳng}:
tách biệt "agent đã ở đâu" khỏi "agent làm được gì"; (P3) \emph{hệ task hai
tầng}: low-level task dẫn dắt khám phá, semantic task cung cấp dữ liệu dài
hạn.

% Không dùng môi trường float `table`: đề cương là phần chèn hành chính, bảng
% của nó không được đánh số chung với bảng các chương, cũng không vào Danh
% sách bảng. (`\caption*` không dùng được vì package `float` phá nó.)
% minipage giữ tiêu đề và bảng trên cùng một trang (center thường vẫn ngắt được)
\begin{center}
\begin{minipage}{\textwidth}
\centering
\textbf{Sáu giai đoạn của pipeline HSGE.}
\vspace{0.6em}

\small
\begin{tabular}{@{}L{3.6cm} L{9.2cm}@{}}
\toprule
\textbf{Giai đoạn} & \textbf{Chức năng chính} \\
\midrule
Biểu diễn graph hai tầng & Tầng thấp: trạng thái trang + cạnh trajectory.
Tầng cao: skill dạng option $\langle I, \pi, \beta \rangle$ + operator có
precondition/effect \\
\addlinespace
Khám phá có định hướng & Hàng đợi ưu tiên frontier (URL đã đi tới so với href
khai thác, có trần theo template); explorer đề xuất low-level task; solver
thực thi trong trình duyệt thật \\
\addlinespace
Phát hiện skill & Cắt đoạn $\rightarrow$ gom cụm trực tuyến $\rightarrow$
trừu tượng hóa thành SkillCard tham số hóa $\rightarrow$ kiểm chứng bằng thực
thi lại trong ngữ cảnh mới \\
\addlinespace
Sinh semantic task & Induction operator nhận biết từ vựng, chuẩn hóa mệnh đề,
cổng chặn tautology; ghép chuỗi effect $\rightarrow$ precondition thành goal
task \\
\addlinespace
Tinh lọc dữ liệu & Judge độc lập chấm lại thành công tự tuyên bố; viết lại suy
nghĩ; dựng dataset dưới ràng buộc prompt parity \\
\addlinespace
Fine-tune SLM & LoRA trên Qwen3.5-9B; phục vụ bằng vLLM; đánh giá bằng bộ chấm
chính thức của WebArena \\
\bottomrule
\end{tabular}
\end{minipage}
\end{center}

\subsubsection*{2.4.2\quad Khám phá có định hướng theo graph}

Frontier được quản lý bằng hàng đợi ưu tiên hai tầng. Tầng 0 là các URL mà
solver thực sự đi tới khi giải task, được nạp thẳng. Tầng 1 là các liên kết
khai thác từ trang, phải qua một trần theo \emph{khuôn} URL (các đường dẫn
chỉ khác nhau ở phần định danh được quy về cùng một khuôn) để những trang
gần trùng nhau không làm ngập frontier. Node có nội dung trùng với node đã
thăm bị bỏ qua; node gây lỗi điều hướng được ghi nhận và bỏ qua thay vì làm
hỏng cả lượt chạy.

Trên mỗi trang, explorer đề xuất các task cụ thể có phân loại (tra cứu thông
tin / điều hướng / thay đổi trạng thái) và ghi nhớ những task đã đề xuất để
tránh lặp. Solver thực thi từng task trong action space cố định gồm click, gõ,
chọn, nhấn phím, điều hướng, thao tác tab, cùng hai kết cục tuyên bố là hoàn
thành và bất khả thi.

\subsubsection*{2.4.3\quad Phát hiện và kiểm chứng skill}

Mỗi trajectory được cắt thành các đoạn con tại các điểm đổi ngữ cảnh, gom cụm
trực tuyến bằng ngưỡng cosine trên embedding, và mỗi cụm lặp lại đủ số lần sẽ
được trừu tượng hóa thành một SkillCard tham số hóa. Skill chỉ được thu nhận
sau khi thực thi lại thành công trong các ngữ cảnh chưa gặp, theo nguyên tắc P1.

\subsubsection*{2.4.4\quad Các độ đo đánh giá}

Độ đo chính là \textbf{tỷ lệ hoàn thành task} (task completion rate) trên bộ
task chính thức của WebArena, chấm bằng bộ đánh giá tự động của benchmark dựa
trên trạng thái cuối, so sánh mô hình fine-tune với chính mô hình nền. Các độ
đo phụ gồm tỷ lệ thành công theo loại task, tỷ lệ episode được judge xác nhận,
và các thống kê về dữ liệu thu thập (số skill, số operator, số cạnh trong
operator graph). Kết quả được lặp lại nhiều lượt để báo cáo trung bình kèm độ
lệch chuẩn và khoảng tin cậy, thay vì một lượt chạy đơn.

\subsubsection*{2.4.5\quad Công nghệ, công cụ và tài nguyên}

\begin{itemize}
  \item \textbf{Ngôn ngữ:} Python $\geq$ 3.11.
  \item \textbf{Trình duyệt:} Playwright + Chromium, qua BrowserGym.
  \item \textbf{LLM vận hành pipeline:} API thương mại, định tuyến theo vai
  trò (explorer, solver, segmenter, abstractor, verifier, operator inducer,
  judge).
  \item \textbf{Huấn luyện và phục vụ:} LLaMA-Factory (LoRA), vLLM.
  \item \textbf{Tài nguyên GPU:} GPU đám mây thuê ngoài, NVIDIA A100-80GB, cho
  fine-tuning và phục vụ suy luận khi đánh giá.
\end{itemize}

\subsection*{2.5\quad Kết quả dự kiến}

\textbf{Mã nguồn mới tạo:}
\begin{itemize}
  \item Pipeline HSGE hoàn chỉnh: khám phá có định hướng theo graph, phát hiện
  và kiểm chứng skill, induction và ghép chuỗi operator.
  \item Bộ công cụ tinh lọc dữ liệu: judge chấm lại episode, bộ viết lại suy
  nghĩ, trình dựng dataset tuân thủ prompt parity.
  \item Bộ script huấn luyện, phục vụ và đánh giá đầu--cuối trên WebArena.
\end{itemize}

\textbf{Dữ liệu và mô hình:}
\begin{itemize}
  \item Tập dữ liệu SFT gồm các trajectory tương tác web có gán mục tiêu, thu
  thập hoàn toàn tự động, kèm siêu dữ liệu về skill và operator.
  \item Checkpoint LoRA của mô hình web agent 9B đã fine-tune.
  \item Bảng so sánh mô hình nền và mô hình fine-tune trên bộ task chính thức
  của bốn website WebArena.
\end{itemize}

\textbf{Tài liệu kỹ thuật:}
\begin{itemize}
  \item Hướng dẫn cài đặt và chạy pipeline từ khám phá đến đánh giá.
  \item Hướng dẫn tái lập thực nghiệm với các lệnh cụ thể.
  \item Báo cáo kinh nghiệm kỹ thuật khi vận hành LLM agent trên website thật.
\end{itemize}

\subsection*{2.6\quad Kế hoạch thực hiện}

\renewcommand{\arraystretch}{1.25}
\begin{longtable}{@{}C{1.0cm} L{5.6cm} C{1.75cm} C{1.75cm} L{2.6cm}@{}}
\caption*{\textbf{Kế hoạch thực hiện khóa luận.}} \\
\toprule
\textbf{STT} & \textbf{Công việc} & \textbf{Bắt đầu} & \textbf{Kết thúc} &
\textbf{Trạng thái} \\
\midrule
\endfirsthead
\toprule
\textbf{STT} & \textbf{Công việc} & \textbf{Bắt đầu} & \textbf{Kết thúc} &
\textbf{Trạng thái} \\
\midrule
\endhead
\bottomrule
\endfoot
1 & Nhận đề tài; khảo sát tài liệu về web agent, khám phá skill và sinh dữ
liệu tự động & 01/2026 & 01/2026 & Hoàn thành \\
\addlinespace
2 & Thiết kế kiến trúc HSGE: graph hai tầng, biểu diễn skill và operator &
02/2026 & 02/2026 & Hoàn thành \\
\addlinespace
3 & Triển khai môi trường WebArena tự host; xây lớp trình duyệt và action
space & 02/2026 & 03/2026 & Hoàn thành \\
\addlinespace
4 & Nộp đơn đăng ký thực hiện đề tài & \multicolumn{2}{C{3.6cm}}{Theo lịch Khoa} &
Hoàn thành \\
\addlinespace
5 & Cài đặt vòng khám phá: quản lý frontier, explorer, solver & 03/2026 &
04/2026 & Hoàn thành \\
\addlinespace
6 & Nộp đề cương cho Khoa & \multicolumn{2}{C{3.6cm}}{Theo lịch Khoa} & Hoàn thành \\
\addlinespace
7 & Cài đặt phát hiện skill: cắt đoạn, gom cụm, trừu tượng hóa, kiểm chứng &
04/2026 & 05/2026 & Hoàn thành \\
\addlinespace
8 & Cài đặt induction operator và ghép chuỗi thành semantic task & 05/2026 &
05/2026 & Hoàn thành \\
\addlinespace
9 & Thu thập dữ liệu trên Reddit và GitLab (12 lượt chạy, 1.174 episode) &
05/2026 & 06/2026 & Hoàn thành \\
\addlinespace
10 & Tinh lọc dữ liệu và dựng tập SFT theo prompt parity & 06/2026 & 06/2026 &
Hoàn thành \\
\addlinespace
11 & Fine-tune Qwen3.5-9B bằng LoRA; dựng hạ tầng phục vụ và đánh giá &
06/2026 & 07/2026 & Hoàn thành \\
\addlinespace
12 & Đánh giá trên bộ task chính thức của GitLab & 07/2026 & 07/2026 & Hoàn
thành \\
\addlinespace
13 & Viết báo cáo khóa luận (Chương 1--3, 5--6, phụ lục) & 06/2026 & 07/2026 &
Hoàn thành \\
\addlinespace
14 & Khắc phục hợp nhất mệnh đề để kích hoạt ghép chuỗi operator; chạy lại thu
thập & 07/2026 & 08/2026 & Đang thực hiện \\
\addlinespace
15 & Đánh giá trên Shopping, Shopping Admin và Reddit & 07/2026 & 08/2026 &
Đang thực hiện \\
\addlinespace
16 & Lặp lại thống kê nhiều lượt và nghiên cứu ablation & 08/2026 & 08/2026 &
Chưa bắt đầu \\
\addlinespace
17 & Hoàn thiện Chương 4 và toàn bộ báo cáo & 08/2026 & 08/2026 & Chưa bắt
đầu \\
\addlinespace
18 & Nộp đơn đăng ký bảo vệ và bảo vệ khóa luận &
\multicolumn{2}{C{3.6cm}}{Theo lịch Khoa (08/2026)} & Chưa bắt đầu \\
\end{longtable}

\vspace{1.5cm}

\noindent
\begin{tabular}{@{}C{6.6cm} C{6.6cm}@{}}
\textbf{XÁC NHẬN CỦA NGƯỜI HƯỚNG DẪN} &
\textbf{TP. Hồ Chí Minh, ngày \ \ \ \ tháng \ \ \ \ năm \ \ \ \ } \\[4pt]
 & \textbf{NHÓM SINH VIÊN THỰC HIỆN} \\[4pt]
\textit{(Ký và ghi rõ họ tên)} & \textit{(Ký và ghi rõ họ tên)} \\
\end{tabular}

\clearpage
